\documentclass[book.tex]{subfiles}
\begin{document}
\chapter{Coarse Graining and Closure Modeling}

\label{chap:coarse_graining_closure}
\noindent\rule{11cm}{0.4pt} \\
\textbf{Prerequisites:} \autoref{chap:qft_basics}  \\
\textbf{Difficulty Level:} ** \\
\noindent\rule{11cm}{0.4pt}
\vspace{5mm}

Automatic coarse-graining methods, based on adaptations of SINDy can be applied to automatically coarse grain equations for turbulence modeling. This approach is capable of automatically discovering symmetries in the original system \cite{bakarji2021data}

\section{A Brief Introduction to SINDy}

Suppose that we have a dynamical system of the form
\begin{align}
    \frac{d}{dt} x(t) &= f(x(t))
\end{align}

where $x(t)$ denotes the state at time $t$ and $f(x(t))$ governs the update dynamics. In many practical applications $f$ may be a sparse function with only a few terms of $x(t)$ necessary to compute the interaction. We can convert this system into a sparse linear regression task by specifying a space of functions, typically polynomials. Let $X$ be the data matrix and $\dot{X}$ its time derivative
\begin{align}
    \Theta(X) &= \begin{bmatrix}
        x_1(t)& \dotsc & x_n(t) & x_1(t)^2 & x_1(t) x_2(t) & \dotsc 
    \end{bmatrix}
\end{align}
We also introduce a sparse set of $k$ coefficients
\begin{align}
    \Xi &= \begin{bmatrix}
        \xi_1 & \dotsc & \xi_k
    \end{bmatrix}
\end{align}
We then need to solve the equation
\begin{align}
    \cdot{X} &= \Theta(X) \Xi
\end{align}
Nonzero coefficients can then be converted back into a symbolic set of equations to gain a simple governing equation for the system from data. We may want to assume that noise is present in which case the equation becomes
\begin{align}
\cdot{X} &= \Theta(X) \Xi + \eta Z
\end{align}
where $Z$ is a Gaussian noise vector. LASSO can be applied to solve this problem.
%"The sparse identification of nonlinear dynamical systems (SINDy) method is applied to probabilistic coarse-graining and closure modeling problems, frequently encountered in turbulence modeling. We show that physical constraints can be encoded in the symmetries of the numerical derivative features operator and conservative probability density function (PDF) equations can be discovered from data. Promising results and challenges are explored. Finally, a short tutorial on discovering a Fokker-Planck equation from Brownian motion simulations of a stochastic differential equation is presented."

%References 

\end{document}